\section{Discussion}
\label{sec:discussion}

CAP-G's value is concentrated and explainable: at triage capacity, weighting the queue by
asset criticality moves mission-critical exposure to the front, nearly doubling
rank-weighted mission precision at no measurable cost to the metric the base scorer was built
for. The honest caveats are equally clear. First, the benefit is \emph{regime-bound}: at
generous capacity it shrinks toward zero, so CAP-G is a triage tool, not a universal
improvement. Second, the benefit is \emph{heterogeneity-bound}: on a uniform fleet it is
mathematically nothing. Third, of the three context dimensions, only criticality
contributed on the mission target as we defined it, a result that argues for simpler
context models and against the assumption that every CMDB attribute helps.

Each of these caveats is worth reading as a result rather than a limitation. The capacity decay
is not a weakness of the scorer but a property of the problem: when the remediation budget is
large enough to cover most high-priority pairs, the order in which they are reached stops
mattering, so any reordering, however principled, converges to the baseline. The value of
asset context is therefore exactly co-located with the scarcity that makes prioritization
necessary in the first place. The heterogeneity bound is equally structural. CAP-G cannot
conjure a distinction the fleet does not contain; on a uniform fleet there is no mission
ordering to recover, and the mechanism control makes this provable rather than merely observed.
That a context-aware method does nothing where context does not vary is the correct behaviour,
and demonstrating it is what separates a real effect from a scoring artifact. The ablation
result, that asset criticality carries the signal while network zone and data sensitivity act as
mild distractors at their pre-registered weights, is the most actionable of the three: it says a
practitioner can deploy a criticality-only context layer and capture essentially all of the
benefit, trading a marginal loss for a large gain in simplicity and explainability. It also
cautions against the common assumption that every attribute already sitting in a configuration
management database is worth feeding into a ranking; more context dimensions are not
automatically better, and on this target two of three were not.

The pre-registered $5$~pp-averaged-over-$K$ bar was, in retrospect, mis-specified for a
quantity that decays with capacity; a triage-regime bar
(e.g., $5$~pp at $K{=}50$) would have been cleanly passed. We leave the bar as locked and
report partial rather than re-cut it post hoc. This is a deliberate methodological stance:
the credibility of a pre-registered study rests on honoring the bar that was fixed in advance,
even when a differently-specified bar would have read more favorably, because the freedom to
choose the threshold after seeing the data is exactly the freedom pre-registration is meant to
remove. Reporting the primary hypothesis as partial, when the effect is real and
interval-separated at every capacity and clears $5$~pp in the triage regime where the tool is
meant to be used, is the price of that discipline and, we argue, an asset to the reader who must
judge the work under adversarial review. For deployment, the operational implication is direct:
agencies running tight remediation windows, the common case, should expect a meaningful
mission-precision gain from a criticality-weighted queue, while agencies with ample capacity
gain little and can prefer the simpler context-blind scorer. The autonomic-computing vision of
systems that continuously assess and reorder their own risk posture~\cite{kephart2003vision}
is the broader setting into which a scorer like this fits, and our results locate where in that
vision asset context adds value and where it saturates.
